\begin{transitionframe}{Entry of female workers to Tokyo's civil service}
    \begin{itemize}
        \item Traditionally male dominated sector.
        \begin{itemize}
        \item Less than 2\% of employees were women.
        
        Concentrated in medical departments. \hyperlink{Department}{\beamergotobutton{Gender ratio by department}}

        \vspace{0.5em}
        
        \item Women were hired as contract workers without tenure. \hyperlink{Occupation}{\beamergotobutton{Gender ratio by Occupation}}

        \end{itemize}
    \
    \pause
        
        \item Tokyo's public sector was occupied by disproportionately high-educated employees.\hypertarget{Descriptive}{}

        \begin{itemize}
            \item 40\% of employees graduated from university. 
            
            $\leftarrow$National average: 3\%.
        \end{itemize}         

        
    \end{itemize}
\end{transitionframe}

\begin{transitionframe}{Japan's Conscription System}
    \begin{wideitemize}
        \item Tokyo civil service began losing employees to the Japanese military starting in 1937.

        \item Reservists were required to leave their civilian positions to serve in the military.

        \textit{Reservists}: Veterans who had previously trained or served.

        \begin{itemize}
            \item Approx. 800 employees left their positions during this period.
            \item High level civil servants were exempt from drafting.
        \end{itemize}

        \item Neither reservists nor employers could not predict when reservists would be recalled to active duty.
    \end{wideitemize}
\end{transitionframe}


\begin{frame}{}
\begin{table}[ht]
\caption{Balance Table of office attributes by drafted status}
\centering
\begin{tabular}{lrrrr}
  \hline
 Year: 1931 & Not Drafted & Drafted & Difference & T‑stat \\ 
  \hline
  Female \%        & 0.02 & 0.02 & 0.00 & -0.33 \\ 
  Civil Servant \% & 0.09 & 0.06 & -0.03 & 1.01 \\ 
  \# of members         & 35.25 & 108.67 & 73.42 & -3.76 \\ 
  \hline 
\end{tabular}
\end{table}

\begin{table}[ht]
\centering
\begin{tabular}{lrrrr}
  \hline
 Year: 1936 & Not Drafted & Drafted & Difference & T‑stat \\ 
  \hline
  Female \%        & 0.04 & 0.02 & -0.02 & 3.07 \\ 
  New Hire \%      & 0.49 & 0.53 & 0.04 & -1.09 \\ 
  Civil Servant \% & 0.06 & 0.09 & 0.03 & -1.41 \\ 
  \# of members    & 18.02 & 56.48 & 38.46 & -4.96 \\ 
  \hline 
\end{tabular}
\end{table}
\note{Here I'm showing a balance table of drafted and non-drafted workers. My main points is that there isn't significant difference in experience with working with women. Both years show small numbers, which underlines how gender balance was skewed in the Tokyo civil service. Notably, I find that larger offices have higher chances of getting drafted. This is concerning factor as office size correlates with drafting and female shares.}
\end{frame}
